% ============================================================================
%  chapter2.tex
%  Registers: From the Register Image Down to the Atom
% ============================================================================

\chapter{Registers: The Bistable Heart of the Machine}

\section{A Snapshot of the Whole Machine}

A process is running. It has just finished one instruction and is
about to start the next when the operating system decides that
another process deserves the CPU. The machine must not lose the
process that was running---a few milliseconds from now it will come
back and finish exactly where it left off.

So the CPU does something extraordinary: it photographs itself. It
copies into memory everything that defines ``where this computation
stands''---which instruction runs next, where the stack points, the
value of every local the compiler is juggling---and later, when the
process returns, it restores the photograph and the computation
continues as if nothing had happened. Chapter~5 will call this
snapshot the \emph{register image}, and Chapter~11 will show a thread
switch saving and restoring it, byte by byte.

\begin{verbatim}
   The register image of an x86-64 process, mid-context-switch:

   RAX  0x000000000000003f     R9   0x00007ffc1004d0
   RCX  0x0000000000000001     R10  0x0000000000000000
   RDX  0x0000000000000000     R11  0x0000000000000246
   RBX  0x0000000000000000     R12  0x0000000000000000
   RSP  0x00007ffc1004f0       R13  0x0000000000000000
   RBP  0x00007ffc1004f0       R14  0x0000000000000000
   RSI  0x00007ffc1004d0       R15  0x0000000000000000
   RDI  0x0000000000000001
   RIP  0x0000000000401167     (the next instruction)
   RFLAGS 0x0000000000000202   (carry, zero, sign, ...)
\end{verbatim}

Sixteen general-purpose registers at eight bytes each, plus the
instruction pointer and the flags: 128 bytes of ordinary, unglamorous
data. And yet those 128 bytes are the whole machine, suspended in
mid-thought. Every computation that will ever run on this CPU spends
its working life inside them.

\begin{osconcept}
None of this saving and restoring happens on its own. The decision to
interrupt a running process, the choice of \emph{which} process runs
next, and the obligation to put the old process's 128 bytes away and
lift another's back out are all made by the \emph{scheduler}, a piece
of the kernel. A thread switch is, in the most literal sense, an
exchange of register images; Chapter~11 will write that exchange by
hand. Registers are hardware, but their contents are handed between
processes only because the OS supervises the transfer.
\end{osconcept}

That is the image this chapter is about. By its end you will know,
bit by bit, what it is made of---because underneath those 128 bytes
lies a ladder of smaller and smaller things. A register. A
flip-flop. Two logic gates. A pair of transistors. And finally, a few
billion silicon atoms, a handful of which have been swapped for
dopants. Let us climb down.

\section{The Register: One Word, Sixty-Four Flip-Flops}
\label{sec:register}

The image is sixteen words. A \emph{word} is the machine's natural
unit (Chapter~1), and each word in the image lives in a
\emph{register}---a small store on the CPU die itself, nanoseconds
away, managed entirely by your program.

A register is not mysterious. It is a row of storage cells, one per
bit: sixty-four cells for a 64-bit register. Each cell is a
\emph{flip-flop}, a circuit that can remember a single 0 or 1, and
all sixty-four flip-flops in a register share one \emph{clock} line
so that the whole word changes at the same instant. When the compiler
emits \code{mov rax, 42}, the CPU drives 42 onto the data lines and
pulses the clock; on that pulse every flip-flop captures its bit, and
\code{rax} now holds 42.

\begin{verbatim}
   A 64-bit register is 64 flip-flops and one clock:

       D0 ----->[D Q]----- Q0
       D1 ----->[D Q]----- Q1
       ...      [D Q]
       D63----->[D Q]----- Q63

                 +-------+
       clock ----> all 64: on the rising edge each captures its D.
                 +-------+
\end{verbatim}

The shared clock is what makes a register \emph{one object} rather
than sixty-four coincidences: the whole word is latched in a single
tick. That is why a register has the byte size of a word, why the
register image is laid out in 8-byte units, and why your program can
treat a register as an indivisible thing.

The sixteen registers are not interchangeable. The machine gives each
one a personality, and the compiler and the hardware agree on it:

\begin{verbatim}
   RAX   accumulator: arithmetic results, function return values
   RCX   counted loops; the first "loop counter" a machine ever had
   RDX   data/remainder; the high half of wide arithmetic
   RBX   a callee-saved scratch register
   RSP   the stack pointer: where the top of the stack is
   RBP   the base/frame pointer: where the current frame starts
   RSI   source index: string ops, second function argument
   RDI   destination index: string ops, first function argument
   R8..R15  a general-purpose pool added by the 64-bit extension
\end{verbatim}

Most of these roles are convention, not law---the CPU will happily do
arithmetic in \code{r15}---but the convention is so deeply baked into
compilers, debuggers, and system call interfaces that your code
should treat it as the layout of the machine. \code{rsp} and
\code{rbp} in particular are \emph{hard} convention: Chapter~6 will
show the stack machinery that depends on them.

\begin{insight}[Why sixteen?]
The number of registers is an accident of instruction encoding. Each
instruction that names two registers must fit those names into a few
bits; \code{R8}--\code{R15} came with a 64-bit extension byte added
to that encoding. A machine with 32 or 128 registers is perfectly
possible---and modern out-of-order CPUs secretly have far more---but
\emph{architecturally}, the number is whatever the instruction format
can address.
\end{insight}

Now the question this chapter descends toward: what is a flip-flop,
that it can hold a bit?

\section{Down One Level: The Flip-Flop}
\label{sec:flipflop}

A flip-flop's ancestor is the \emph{SR latch}, and the SR latch is
almost absurdly simple: two logic gates wired so that each gate's
output is one of the other gate's inputs. With no inputs asserted,
the circuit is just a loop. And that loop is a memory.

\begin{verbatim}
   The SR latch: two NOR gates, cross-coupled.

         S -------------------------+
                                     v
                        +-----------------------+
                        |          NOR          |
                        |                       |----------- Q
                        +-----------------------+
                              ^
                              |   (Qbar feeds back)
                        +-----------------------+
                        |          NOR          |
                        |                       |----------- Qbar
                        +-----------------------+
                              ^
         R --------------------+
\end{verbatim}

This is the drawing you will find in Wikipedia's flip-flop article
and in every digital-logic textbook; in vacuum-tube form it dates to
the 1918 Eccles--Jordan trigger relay. A NOR gate outputs 1 exactly
when both its inputs are 0, which gives the latch a tiny truth table:

\begin{verbatim}
   S   R   action
   --- --- -----------------------------------
   1   0   set     : Q = 1 (the latch remembers 1)
   0   1   reset   : Q = 0 (the latch remembers 0)
   0   0   hold    : Q stays wherever it was
   1   1   forbidden: both outputs collapse; never drive this
\end{verbatim}

Force S to 1 and Q rises; the lower NOR sees that high and its output
falls; and when you let go of S, the loop is on its own---and it
stays exactly where it is, because each gate holds the other in
place. That is the entire mechanism of computer memory in one
sentence: \emph{a loop that reinforces itself}.

A bare latch is transparent---a stray signal can flip it at any time.
A CPU wants its registers to change only at agreed instants, so the
flip-flops it uses are \emph{edge-triggered D flip-flops}: one data
input, one clock, and the rule that the output changes only on the
clock's rising edge, and only to whatever the data input held at that
exact moment. A D flip-flop is two latches in series---master and
slave, opened on opposite clock phases---so it samples, then holds,
with no window for accidental change. Sixty-four of them sharing one
clock is the register of Section~\ref{sec:register}.

\begin{verbatim}
   A D flip-flop as master and slave, one data path:

         D ---->[master latch]----[slave latch]----> Q
                       |                  |
                       +--------+---------+
                                |
                                v
                          clock (inverted for one of them)
   On the rising edge the master opens and samples D; the slave is
   locked. On the falling edge the master locks and the slave opens,
   copying the sampled value to Q. The output therefore changes only
   around the clock edge, never in between.
\end{verbatim}

\section{Why the Loop Holds: Bistability and the Math of Feedback}
\label{sec:bistable}

The latch holds because it is \emph{bistable}: it has exactly two
stable states, and every other state is shoved toward one of them.
The mathematics is clean, and worth holding in your hands.

Model each NOR gate as a sharp inverting element. With $Q$ and
$\bar{Q}$ the two stored nodes, the pair of gates imposes

\begin{equation}
Q = \overline{S \lor \bar{Q}}, \qquad \bar{Q} = \overline{R \lor Q},
\end{equation}

where $\lor$ is logical OR and the overbar is negation. In the hold
condition ($S = R = 0$) this collapses to

\begin{equation}
Q = \overline{\bar{Q}}, \qquad \bar{Q} = \overline{Q},
\end{equation}

which has exactly two solutions: $Q=1,\ \bar{Q}=0$ and $Q=0,\
\bar{Q}=1$. A third, formal solution---$Q$ and $\bar{Q}$ both
halfway---exists on paper but not in practice: it is an
\emph{unstable} equilibrium. The tiniest noise tips the loop, and
positive feedback then races it all the way to one of the stable
corners. That razor's edge is called \emph{metastability}, and it is
why real hardware occasionally hiccups when two signals race.

The analogue view says the same thing with a picture. An inverter's
output $v_{out}$ is a function $f$ of its input; two inverters in a
loop must satisfy $v_{out} = f\bigl(f(v_{out})\bigr)$. Because $f$
is inverting and steeper than $-1$ at its midpoint, the composition
crosses the diagonal three times: the two stable points and the
unstable middle.

\begin{verbatim}
   v_out
   1 |  x                     stable: Q = 1
     |    .
     |     .     the middle diagonal crossing is the
     |       .        .       metastable point: unstable
     |        .    .
     |           .  x         stable: Q = 0
   0 +----------------------------- v_in
\end{verbatim}

That is the mathematical reason a register does not ``half-remember'':
the loop has no resting place between 0 and 1, so it cannot rest
there. A flip-flop is \emph{static} memory---it holds its two stable
voltages as long as power flows, needing no help. There is a second,
dynamic family, the DRAM capacitor, whose retention is a race against
leakage and whose math is even more explicit. Both families are built
from the same two ingredients, which is the next level down.

\section{The Clock: Making a Million Flip-Flops Agree}
\label{sec:clock}

Every register in the CPU is edge-triggered, and edge-triggering only
means something if every edge arrives at the same instant across the
whole die. That agreement is the job of the \emph{clock}: a single
square wave, distributed to every flip-flop, that ticks once per
cycle. At roughly 3--5 billion ticks a second in a modern CPU, each
tick lasts about a third of a nanosecond.

\begin{verbatim}
   The clock, drawn against time (t is running to the right):

   clock   ____|----|____|----|____|----|____|----|____
                  ^         ^         ^         ^
              rising edges: on each one, every flip-flop in the
              CPU samples its D input and changes exactly then.
\end{verbatim}

A synchronous machine is a heartbeat: flip-flops remember, and the
combinational logic \emph{between} registers---gates wired in a
straight line, no loops---does the computing during the quiet part of
each cycle. One tick to sample, one interval to compute, one tick to
latch the answer, and the whole machine advances in lockstep. This is
why your program sees registers change ``all at once'' and never
``in mid-thought'': the clock simply does not allow half-updated
words.

\subsection{Setup and Hold: The Flip-Flop's Polite Window}

A flip-flop samples on the edge, but it needs a moment of peace around
that edge. The data must be settled and steady for a short time
\emph{before} the edge---the \emph{setup time}, $t_{su}$---and must
stay put for a short time \emph{after} it---the \emph{hold time},
$t_{h}$. These are physical facts about the transistors inside the
flip-flop, on the order of tens of picoseconds.

\begin{verbatim}
   The window a flip-flop demands around a rising edge:

   clock    ____|----|________|----|____
                        ^
   D input  ____|_______________|________
             <-- stable high --><-- stable low
                        |<- t_su ->|<- t_h ->|
                        D must not change inside this window.

   Violate it, and the flip-flop does not guarantee a clean 0 or 1:
   it may sit halfway, decide randomly, or take a while to decide.
\end{verbatim}

Setup and hold are the price of agreeing on \emph{when}. The clock
must arrive at every flip-flop at nearly the same time (a skew of a
few dozen picoseconds is acceptable), the combinational logic must
finish within the cycle, and no signal may sneak through a latch
while its data is still being written. Engineers design for these
rules with static timing analysis; your program never sees them,
because the CPU is built to meet its own deadlines by construction.

\subsection{Metastability in the Wild: Crossing Asynchronous Borders}

Metastability is not an academic footnote. It appears the moment
something \emph{asynchronous} touches the clocked machine: a keyboard
key, a network packet, a disk interrupt. Those signals do not march to
the CPU's clock, so when a flip-flop samples them, the sample can land
right inside the forbidden window---and the flip-flop goes metastable,
hesitating between 0 and 1 for a while before tumbling to one side.

The trick is to not care which side it falls on. A
\emph{synchronizer} is two flip-flops in a row fed by the same clock:
the first may go metastable, but it has a whole clock cycle to settle,
and the second flip-flop samples it \emph{after} that, by which time
it is a solid 0 or 1. We lose a cycle of latency; we gain a guarantee
that the value the logic sees is a clean bit.

\begin{verbatim}
   Two flip-flops turning an async input into a safe signal:

        async in --->[D Q]---->[D Q]----> into the clocked logic
                      flop 1    flop 2
                         |         |
                         +----+----+
                              |
                              v
                          the one clock (both flops)

   flop 1 has one full clock to resolve; flop 2 then samples a
   settled value. The value is "arbitrary," but it is never
   mid-flight, and that is what the rest of the machine needs.
\end{verbatim}

That is why hardware ``occasionally hiccups'' in the face of racing
signals: no synchronizer is perfect, and with enough samples the
first flop occasionally takes longer than a cycle to decide. The
failure rate is not zero---it is merely small---and how small is a
measured engineering quantity. It is also exactly why every external
interrupt in the kernel arrives through hardware that has already
been synchronized before it ever reaches a register.

\begin{researchnote}[Metastability is measured, not guessed]
The frightening part of metastability is that a synchronizer can
theoretically sit unresolved forever. In practice it cannot, but the
failure rate is governed by a stark formula. For a two-flip-flop
synchronizer given a resolution time $t_r$, the mean time between
failures is roughly

\begin{equation}
\text{MTBF} \approx \frac{e^{t_r / \tau}}{T_0\, f_\text{clk}\, f_\text{data}},
\end{equation}

where $f_\text{clk}$ and $f_\text{data}$ are the clock and data
frequencies and $T_0$ and $\tau$ are constants of the flip-flop
technology. The exponential in the numerator is the whole story:
each extra clock cycle of resolution time buys an astronomically
larger MTBF. Thomas Chaney and Charles Molnar proved the phenomenon
was fundamental, not a manufacturing flaw, in 1973, and the formula
appears in nearly every serious digital-design text since.

\paperref{T.~J.~Chaney, C.~E.~Molnar, ``Anomalous Behavior of
Synchronizer and Arbiter Circuits,'' IEEE Transactions on Computers,
1973}{https://ieeexplore.ieee.org/document/1672263}
\end{researchnote}

\begin{ruleofthumb}[Synchronize before you trust]
In software terms, treat every value that arrives from outside the
clocked core---a pressed key, a network byte, a changed file---as
asynchronous and therefore suspect. The hardware synchronizer turns a
``maybe metastable'' bit into a ``settled, but arbitrary'' bit; your
program's job is to re-validate that arbitrary value (parse it, check
it) before trusting it. The register image itself never needs this
treatment: everything in it lives on the clock.
\end{ruleofthumb}

\section{Down Another Level: Transistors and Capacitors}
\label{sec:transistors}

The gates and inverters above are not symbols; each is built from
\emph{field-effect transistors}, which are voltage-controlled
switches. A voltage on the gate opens or closes a channel for
current. A CMOS inverter is two transistors: an n-type one that
conducts when its input is high and pulls the output low, and a
p-type one that conducts when the input is low and pulls the output
high.

\begin{verbatim}
   A CMOS inverter: one p-type and one n-type transistor.

           Vdd
            |
           [P]    p-type: on when input is 0
            |----- output
           [N]    n-type: on when input is 1
            |
          GND

   input = 1 -> N on, P off -> output 0
   input = 0 -> P on, N off -> output 1
\end{verbatim}

A static RAM (SRAM) cell---the flip-flop of Section~\ref{sec:flipflop}
drawn in silicon---is six transistors: two cross-coupled inverters
(four transistors) that hold the bit, plus two \emph{access}
transistors that connect the cell to the bit lines when the row's
word line is high. Read or write, the cell is a tiny latch, exactly
the loop of the previous section.

\begin{verbatim}
   The 6-transistor SRAM cell: the two inverters, plus access.

            Vdd                     Vdd
             |                       |
            [P1]                    [P2]     <- pull-up transistors
             |                       |
             +--- Q ----+---- Qbar --+
             |          |            |
            [N1]       [A1]        [A2]     <- access transistors
             |          |            |
           GND     bit line     bit line bar
                         \        /
                          +------+--------- word line (their gates)

   word line low  : A1, A2 open   -> the cell floats, holding Q
   word line high : A1, A2 closed -> bit lines read or write Q
   The two inverters actively pump Q and Qbar to the rails, which is
   why SRAM never needs refreshing.
\end{verbatim}

A dynamic RAM (DRAM) cell is leaner: one access transistor and one
\emph{capacitor}. The bit is simply whether the capacitor holds
charge. That frugality is why DRAM is cheap and dense---and its cost
is that the charge leaks away and must be \emph{refreshed} on a
timer.

\begin{verbatim}
   A DRAM cell: one transistor, one capacitor.

        word line
            |
            v
           [T]     the access transistor
            |
   bit line ---+====   the storage capacitor (the bit)
               |
              gnd
\end{verbatim}

Here is the retention math, concrete and vivid. A capacitor stores
charge $Q = C\,V$. A leakage current $I$ drains it at

\begin{equation}
\frac{dV}{dt} = -\frac{I}{C},
\end{equation}

and if the leak behaves like a resistance $R$, the voltage decays as

\begin{equation}
V(t) = V_0\,e^{-t/RC}.
\end{equation}

Take a modern cell: roughly $C \approx 25$ fF and a standby leakage
of about a picoampere. Then $dV/dt \approx 40$ V/ms, so a stored
voltage of about a volt falls below the sensing threshold within ten
to twenty milliseconds. That is why the JEDEC standard insists every
DRAM row be refreshed at least every 64 ms---the memory must be
re-read and topped up before its bits fade into noise. The register's
flip-flops need no such clock: their inverters actively pump the
stored node back to the rail, which is why the register image holds
its values while the power is on.

\begin{verbatim}
   Two families, one difference:

                         holds the bit        refreshed?
   register flip-flop    active inverter      never (static)
   SRAM cell             active inverter      never (static)
   DRAM cell             a passive capacitor  every ~64 ms

   Static families actively pump their bit; dynamic families store
   charge and fight leakage. Everything else is the same loop.
\end{verbatim}

\begin{insight}[The cache hierarchy is this section in disguise]
Chapter~1's memory hierarchy is a ladder of exactly these cells.
L1/L2/L3 cache is SRAM---cross-coupled inverters, no refresh, fast
but expensive. Main memory is DRAM---one transistor and a capacitor,
cheap and dense, refreshed continuously by the memory controller.
And the register file is the fastest SRAM of all, right at the top
of the ladder. Every ``level'' you met in Chapter~1 is one of the two
circuits drawn here, scaled up.
\end{insight}

\section{The Bottom: Chemistry, Silicon, and a Handful of Electrons}
\label{sec:chemistry}

What are those transistors and capacitors made of? Silicon, arranged
in a crystal, with a few of its atoms deliberately replaced. This is
the chemistry underneath every bit your program has ever stored.

Pure silicon is a semiconductor. Each atom shares four valence
electrons in a perfect lattice, so at room temperature almost nothing
conducts. \emph{Doping} changes that by substituting a handful of
atoms:

\begin{itemize}
  \item \emph{Phosphorus} brings five valence electrons and donates
    a free one: the material becomes \emph{n-type}, full of mobile
    electrons.
  \item \emph{Boron} brings three valence electrons and leaves a
    missing bond---a \emph{hole} that moves like a positive
    charge---and the material becomes \emph{p-type}.
\end{itemize}

A transistor is a p-type body with n-type source and drain, a few
nanometres of silicon dioxide on top, and a metal gate. Raise the
gate voltage and its electric field pulls electrons to the surface,
forming an \emph{inversion channel} that connects source to drain.
No field, no channel, no current: that is the voltage-controlled
switch of the previous section, made of silicon and a sprinkling of
boron and phosphorus atoms. At doping levels near $10^{17}$ atoms per
cubic centimetre, roughly one silicon atom in a million is a dopant;
the ``free'' electrons that carry a signal are donations from those
rare, chosen atoms.

\begin{verbatim}
   The chemistry of a switch:

        gate
          |
         SiO2    (a nanometre of glass)
     n+  | p-type | n+
     S   | channel| D
         |  (electrons pulled in by the gate's field)

     electrons are the carriers; holes are their absence.
\end{verbatim}

Leakage, the enemy of the DRAM capacitor, is a chemistry story too.
An ``off'' transistor still lets a few carriers slip through the
barrier, and subthreshold leakage grows exponentially as devices
shrink. A modern DRAM cell stores its bit as only tens of thousands
of electrons; every time the cell shrinks it holds fewer, leaks
faster, and must be refreshed more aggressively. The famous scaling
rule---R.~H.~Dennard's 1974 prediction that smaller transistors would
stay fast and cool---broke precisely because of this chemistry of
tunnelling and leakage. That is why the industry moved to new
dielectrics, finFETs, and three-dimensional transistors: the deepest
detail on this ladder is where Moore's law meets its limit.

\begin{verbatim}
   The full descent, one ladder:

   register  -> 64 flip-flops sharing a clock
   flip-flop -> two cross-coupled inverting gates, edge-triggered
   gate      -> a pair of transistors (CMOS inverter, NOR, ...)
   transistor-> a sliver of doped silicon with a field-effect switch
   silicon   -> a crystal of atoms, a handful doped with P or B
   bit       -> the presence or absence of electrons at one node
\end{verbatim}

\section{Back Up: The Image, Understood}
\label{sec:backup}

Let us return to the opening scene. The process is interrupted; the
CPU photographs itself; 128 bytes of register image go into memory.
Now you know what those bytes are:

\begin{verbatim}
   one bit      = one flip-flop = two cross-coupled gates
                = four to six transistors = a loop of doped silicon
   one register = 64 flip-flops sharing a clock   (8 bytes)
   the image    = 16 registers + RIP + RFLAGS     (128 bytes)
\end{verbatim}

The register image is not a metaphor. It is 128 bytes of bistable
loops---each loop a pair of gates, each gate a pair of transistors,
each transistor a sliver of doped silicon---sitting on the CPU die
and holding the exact state of a computation because each loop
actively refuses to change. That refusal, powered by positive
feedback and paid for elsewhere with leakage and refresh, is the
deepest reason the data structures, processes, and databases of this
book work at all.

\begin{insight}[Every byte you store is this story]
The 128-byte register image is the fastest memory there is, and its
cousins---L1 cache, SRAM, DRAM, disk---are the same story at larger
scale. Cache is SRAM: more cross-coupled inverters, still static,
still loop-held. RAM is DRAM: capacitors racing leakage, refreshed
64 times a second. The hierarchy of Chapter~1 is, from the bottom up,
a ladder of the very same physics you just descended.
\end{insight}

\section{Registers at Work: The Datapath and the ALU}
\label{sec:datapath}

The descent is complete, so let us watch the image \emph{work}. A
register's job is not to sit still; it is to feed a computation. The
machine that does the feeding is the \emph{datapath}: the register
file (all sixteen registers in one bank), the arithmetic/logic unit
(the \emph{ALU}), and the wires that connect them. Every instruction
you will ever compile is a script for this small, closed system.

\begin{verbatim}
   The instruction loop of one add rax, rbx, step by step:

        +--------+      +--------+      +--------+
        |  RIP   |      |  RAX   |      |  RBX   |   register file
        +--------+      +--------+      +--------+
           |    ^           ^               ^
           |    |           |               |
           |    |       read port A     read port B
           |    |
           |  +-----------------------------+
           |  |          control logic      |  "add"
           |  +-----------------------------+
           |            |          |
           |            v          v
           |        +---------------+
           |        |     ALU       |  computes RAX + RBX
           |        +---------------+
           |            |
           |      +-----+------------+-----------------+
           |      |                  |                 |
           |      v                  v                 v
           |  write port -> RAX   RFLAGS           result discarded
           |                          (CF, ZF, SF, OF set)
           +--- RIP advances past the instruction; repeat forever
\end{verbatim}

1. \textbf{Fetch.} The instruction pointer \code{rip} holds the
address of the next instruction's bytes in memory. The CPU reads
those bytes and advances \code{rip} past them.
2. \textbf{Decode.} The control logic works out what the bytes mean:
``add the value in \code{rbx} to the value in \code{rax}.''
3. \textbf{Read.} The register file places \code{rax} on read port
A and \code{rbx} on read port B. Two reads in the same tick is what
the two ports are \emph{for}.
4. \textbf{Execute.} The ALU computes the sum and, as a side effect,
latches the result's meaning into \code{rflags}: whether it was zero
(\code{ZF}), negative (\code{SF}), carried out (\code{CF}), or
overflowed (\code{OF}).
5. \textbf{Write.} On the clock edge, the result is written back
into \code{rax} through the write port, and the whole loop repeats.

Three registers you already know appear here in their true roles.
\code{rip} is the machine's cursor into the program; \code{rflags}
is the machine's memory of what the last computation \emph{meant};
and the register file is the working surface the ALU reads from and
writes to. An add is a fetch, a decode, a read, an add, a write---and
a handful of clock cycles, in a machine that does billions of them a
second.

\begin{insight}[Architectural vs.\ physical registers]
The sixteen registers your program sees are the \emph{architectural}
register file: the ones Chapter~5 saves and Chapter~11 restores. A
modern out-of-order CPU keeps a much larger \emph{physical} register
file underneath---hundreds of cells---and renames the architectural
names onto physical cells so that independent instructions can use
the same \code{rax} in the same tick without colliding. Your program
only ever sees the sixteen names; the renaming is invisible magic
that makes the datapath faster. The image you save is still 128
bytes, because the contract with your program is the architectural
file.
\end{insight}

\section{RFLAGS, Bit by Bit: How a Comparison Decides a Branch}
\label{sec:rflags}

The flags are easy to ignore and impossible to avoid. Every
arithmetic instruction sets them; every conditional branch reads
them; and compilers generate code that lives and dies by them. In the
register image they are a single 64-bit word, but only a handful of
the bits carry meaning:

\begin{verbatim}
   The meaningful bits of RFLAGS (x86-64):

   bit   name  meaning
   ---   ----  ------------------------------------------------
   0     CF    carry: unsigned overflow from the last operation
   2     PF    parity: the low byte of the result had even parity
   4     AF    auxiliary carry: bit 3 carried (BCD arithmetic)
   6     ZF    zero: the result was exactly zero
   7     SF    sign: the result is negative (its top bit is set)
   8     TF    trap: single-step, used by debuggers
   9     IF    interrupt enable: maskable interrupts are allowed
   10    DF    direction: string instructions move up or down
   11    OF    overflow: signed overflow from the last operation
\end{verbatim}

The six arithmetic flags---\code{CF, PF, AF, ZF, SF, OF}---are the
ALU reporting on its own work. The trick that makes them powerful is
that a comparison is just a subtraction that throws away its result.
The instruction \code{cmp a, b} computes $a - b$ and sets the flags
exactly as a subtraction would, but stores nothing. The branch that
follows then reads the flags as an answer.

Consider the simplest possible overflow, on paper:

\begin{verbatim}
   add rax, rbx      rax = 0x7fff...ffff   rbx = 1

   result:           0x8000...0000
     unsigned view:  2^63 - 1 + 1 = 2^63, which fits  -> CF = 0
     signed view:    (2^63 - 1) + 1 = 2^63, which as a signed
                     64-bit number is negative -> overflow -> OF = 1
     zero?  no -> ZF = 0       sign? top bit set -> SF = 1
\end{verbatim}

One operation, two honest answers: \code{CF} reports the truth to
unsigned arithmetic, \code{OF} to signed. That is why there are two
families of conditional jumps. Unsigned comparisons (\code{jb}, \code{ja},
``below/above'') test \code{CF} and \code{ZF}; signed comparisons
(\code{jl}, \code{jg}, ``less/greater'') test \code{SF} against
\code{OF}. The same bytes, the same flags, two different readings.

Let a compiler show it. This C function

\begin{lstlisting}[caption={A branch the compiler will turn into flags.}]
int max(int a, int b) {
    if (a > b) {
        return a;
    } else {
        return b;
    }
}
\end{lstlisting}

compiles (unoptimized, so the branch is visible) into something very
close to this:

\begin{verbatim}
   max:
       pushq  %rbp
       movq   %rsp, %rbp
       movl   %edi, -8(%rbp)     a lives on the stack
       movl   %esi, -12(%rbp)    b lives on the stack
       movl   -8(%rbp), %eax     eax = a
       cmpl   -12(%rbp), %eax    compute a - b, set flags, keep nothing
       jle    .L2                if a <= b (signed) go load b
       movl   -8(%rbp), %eax     the a > b path: return a
       jmp    .L3
   .L2:
       movl   -12(%rbp), %eax    the else path: return b
   .L3:
       popq   %rbp
       ret
\end{verbatim}

(Exact output varies with compiler and version, but the shape is
universal.) The compare sets \code{CF}, \code{ZF}, \code{SF}, and
\code{OF} from $a - b$; \code{jle} then jumps if the subtraction said
``less than or equal'' in signed terms---that is, if $\text{ZF}=1$ or
$\text{SF} \ne \text{OF}$. A modern optimizer notices the branch is
short and replaces the whole dance with \code{cmov}, a conditional
\emph{move} that reads the same flags without branching. The flags
are the machine's internal vocabulary for ``how did that
comparison go?'', spoken in every branch your program ever takes.

\begin{insight}[Flags are a side effect]
Arithmetic sets the flags whether you want them or not. A stray
\code{add} in unrelated code silently rewrites \code{ZF}, which is
why the hardware convention is that branches must immediately follow
the comparison they depend on---and why compilers reorder your code
with such care. In assembly, a comparison is not data until the
branch consumes it. This pairing---set flags, consume flags---is the
lowest-level control flow your code will ever see.
\end{insight}

\begin{ruleofthumb}[Unsigned or signed: pick by the type]
When you compare two \code{size\_t} or \code{unsigned} values, the
compiler emits unsigned branches testing \code{CF} and \code{ZF};
when you compare \code{int} or \code{long}, it emits signed branches
testing \code{SF} against \code{OF}. Mixing the two is the classic
one-off bug (``\code{size\_t} is huge when it underflows'') that
Chapter~1's Loop Counter Rule exists to prevent. The flags remember
both readings; the instruction chooses which to consult.
\end{ruleofthumb}

\section{Practical Exercise: See Your Program's Register Image}
\label{sec:practical}

The opening scene of this chapter was a register image, but you have
probably never seen your \emph{own} program's registers. Let us fix
that. Two tools show you the fastest memory there is: a debugger to
dump it live, and the compiler to show how your C becomes register
moves.

\subsection{Dumping the image with gdb}

Write a tiny program, compile it with debug symbols, and stop it
inside a function:

\begin{lstlisting}[caption={max.c: the function from the previous section.}]
#include <stdio.h>

int max(int a, int b) {
    if (a > b) {
        return a;
    } else {
        return b;
    }
}

int main(void) {
    printf("max(7, 11) = %d\n", max(7, 11));
    return 0;
}
\end{lstlisting}

Compile with debug info and run it under the debugger:

\begin{verbatim}
   $ gcc -O0 -g -o max max.c
   $ gdb ./max
   (gdb) break max
   (gdb) run
   (gdb) info registers
\end{verbatim}

Because we break at the \emph{entry} of \code{max}, before any of its
own instructions run, the general-purpose registers still show the
calling convention in action: the first two integer arguments arrived
in \code{rdi} and \code{rsi}. You are looking at the register image
of a live computation---the 128 bytes of Chapter~2, mid-flight.

\begin{verbatim}
   rax            0x0                 0
   rbx            0x0                 0
   rdi            0x7                 7        <- first argument
   rsi            0xb                 11       <- second argument
   rsp            0x7fffffffe3d8      0x7fffffffe3d8
   rbp            0x7fffffffe3d0      0x7fffffffe3d0
   rip            0x401136            0x401136 <max>
   eflags         0x202               [ IF ]
\end{verbatim}

Try \code{next} a few times and watch \code{eflags} change as the
\code{cmp} and \code{jle} execute; then \code{continue} and watch the
return value land in \code{rax}. This is not a simulation. Every value
you see is held by cross-coupled loops of doped silicon, and a
\code{context switch} would freeze exactly these bytes to disk-shaped
memory and restore them later.

\subsection{Following your data through the compiler}

The debugger shows you the live image; the compiler shows you how it
was produced. Ask \code{gcc} for assembly instead of a binary:

\begin{verbatim}
   $ gcc -O2 -S max.c
   $ cat max.s
\end{verbatim}

At \code{-O2} the optimizer has made the branch disappear---the
\code{cmov} mentioned in the previous section---and the whole
function is a handful of instructions. The calling convention is
readable in the input names: \code{edi} and \code{esi} (the low 32
bits of \code{rdi} and \code{rsi}) enter, and the answer leaves in
\code{eax}. If you know where the arguments and the return value
live, you can read almost any function's skeleton without a manual.

\begin{verbatim}
   The calling convention (x86-64 System V), in one line:

   args 1-6 : rdi, rsi, rdx, rcx, r8, r9    return: rax
   callee-saved: rbx, rbp, r12-r15          caller-saved: the rest
\end{verbatim}

This convention is part of why the register image has the shape it
does. The callee-saved registers are the ones a function must restore
before returning---so a context switch only needs to save \emph{its}
working set, while the hardware and the compiler split the rest by
this rule. When Chapter~5 serializes a register image and Chapter~11
saves it across a thread switch, this table is the key to what gets
saved, and when.

\subsection{What to do when registers are not enough}

A running program needs more than sixteen words. Function arguments
beyond the sixth, large locals, and array indices all spill into
memory---onto the stack and heap of Chapter~1. The register file is
the fastest memory a program has, but it is also the smallest, and
the compiler's endless job is to decide which words deserve a
register and which must be spilled. Run \code{gcc -O0 -S} on the same
\code{max.c} and you will see the spilled version from the previous
section: everything parked on the stack, registers barely used. Run
\code{-O2} and the spills vanish. That trade---registers for speed,
memory for capacity---is the oldest trade in computing, and it is
exactly the hierarchy of Chapter~1 playing out inside a single
function.

\begin{ruleofthumb}[The register keyword is a fossil]
C has a keyword, \code{register}, that asks the compiler to keep a
variable in a register. Modern compilers ignore it almost entirely;
they have better information than you do about which words are hot,
and they reallocate freely. The reason it exists is that C was
written for machines where the answer was not obvious and the
compilers were not good. Let your compiler do its job, and reserve
your attention for the part it cannot decide: what \emph{lifetime}
your data needs, which is the stack-versus-heap question of
Chapter~1.
\end{ruleofthumb}

\section{Structures in Bytes}
\label{sec:structures}

The register file on x86-64, given byte addresses (each register is a
word, 8 bytes):

\begin{verbatim}
   general-purpose registers (16 x 8 bytes = 128 bytes):

   RAX   0     RCX   8     RDX  16     RBX  24
   RSP  32     RBP  40     RSI  48     RDI  56
   R8   64     R9   72     R10  80     R11  88
   R12  96     R13  104    R14  112    R15  120

   special registers:
   RIP      the instruction pointer: which program byte runs next
   RFLAGS   status bits: carry, zero, sign, overflow, direction...

   one D flip-flop: 1 bit
   one register:    64 flip-flops = 8 bytes
   the image:       128 bytes of the fastest storage that exists
\end{verbatim}

RSP is simultaneously the stack pointer and one of the sixteen
general registers; RIP and RFLAGS sit just outside the 128-byte file
but are saved with it on every context switch. This byte map is
exactly what Chapter~11's thread-scheduling code copies out and back
in, and it is the ``$\sim$100s B, $\sim$1 ns'' row of Chapter~1's
hierarchy table, made visible.

Each register is itself a small structure: one storage unit with
several overlapping views. The \code{rax} word can be addressed as a
full 64 bits, its low 32, its low 16, or its individual bytes:

\begin{verbatim}
   RAX, seen as overlapping windows (all of them are the same bytes):

   byte 0   1   2   3   4   5   6   7
   +---+---+---+---+---+---+---+---+
   | AL| AH|   AX (16 bits)         |
   |   EAX (32 bits)                |
   |   RAX (64 bits)                |
   +---+---+---+---+---+---+---+---+

   AL   = low 8 bits     AH  = bits 8..15
   AX   = low 16 bits    EAX = low 32 bits    RAX = all 64

   The register file's second word layer: RAX at bytes 0..7, RCX at
   8..15, and so on down the file. Writing EAX zeroes the top 32
   bits of RAX; writing AL leaves them alone.
\end{verbatim}

This partial-register trick is why the register image in the opening
scene is a clean 128 bytes while the machine inside happily reads
half-words and bytes. The file is 128 bytes of structure---sixteen
8-byte words---and every word is itself a byte map.

The flags word is a structure too, and a bit-level one: each status
flag is a single bit at a fixed position, which is why the kernel
reads and writes \code{RFLAGS} with bit tests and masks rather than
as a number. Section~\ref{sec:rflags} listed the positions; here they
are as bytes and bits:

\begin{verbatim}
   RFLAGS as a byte map (meaningful bits only):

   +-+-+-+-+-+-+-+-+-+-+-+-+-+-+-+-+-+-+-+-+-+-+-+-+-+-+-+-+
   |...|OF|DF|IF|TF|SF|ZF| AF |PF| |CF|
   +-+-+-+-+-+-+-+-+-+-+-+-+-+-+-+-+-+-+-+-+-+-+-+-+-+-+-+-+
      11 10  9  8  7  6   4  2    0

   the other bits are reserved and stay 0, which is why a clean
   flags word reads 0x202 (IF, plus the always-on reserved bits).
\end{verbatim}

Two more banks live beside the 128-byte file and are saved with it on
context switches: the \emph{vector registers} \code{XMM0}--\code{XMM15}
(sixteen 128-bit words, used by floating point and SIMD) and the
\emph{segment registers} \code{CS}, \code{DS}, \code{ES}, \code{FS},
\code{GS}, \code{SS}. Most programs touch only the general file, but
any code that does heavy float or string work will find its state
spread across all three banks. When you read a register dump, these
are the three blocks you are looking at.

\section{Summary}

\begin{verbatim}
   The register model, complete:

   register     a word of storage on the CPU die
   flip-flop    one bit: two cross-coupled inverting gates
   SR latch     set / reset / hold; the ancestor of all of it (1918)
   D flip-flop  captures D on a clock edge; 64 of them = a register
   bistability  two stable states; metastable only on the razor edge
   the clock    the shared edge that makes the whole machine agree
   setup/hold   a flop needs quiet around the edge; break it -> meta
   synchronizer two flops that absorb an async signal's uncertainty
   SRAM         6 transistors, static, no refresh
   DRAM         1 transistor + 1 capacitor, leaking, refreshed
   retention    V(t) = V0 e^(-t/RC): why DRAM refreshes every 64 ms
   datapath     RIP -> decode -> register file -> ALU -> RFLAGS
   RFLAGS       CF PF AF ZF SF OF: the ALU reporting on its work
   cmp/jcc      subtract, keep only the flags, branch on the flags
   chemistry    silicon + phosphorus/boron = the switches

   Rules that keep the model straight:
     1. a register holds its value only while power flows;
     2. writes happen on clock edges, all bits of a word together;
     3. DRAM must be refreshed; SRAM and registers must not;
     4. the register image is a byte map: 16 x 8 = 128 bytes;
     5. async inputs get synchronized; the register file never does;
     6. comparisons set flags; branches read them; never separate them.
\end{verbatim}

\begin{researchnote}[Two century-spanning inventions]
The flip-flop is older than the transistor. William Eccles and
F.~W.~Jordan patented their ``trigger relay''---two amplifying valves
cross-coupled in a positive feedback loop, the same loop as
Section~\ref{sec:flipflop}---in 1918, and a quarter-century later it
was storing the state of the Colossus codebreaking computers. The
gates-drawn diagram of the SR NOR latch in this chapter is the direct
descendant of their two vacuum tubes.

The other half of the story is the scaling rule that made registers
and DRAM cheap. R.~H.~Dennard and colleagues showed in 1974 that
shrinking transistors could keep power density constant while running
faster; the chemistry of leakage in Section~\ref{sec:chemistry} is
the reason that law eventually broke, and it is the same physics that
forces DRAM to refresh today. The metastability side of the story
has its own founding paper: Chaney and Molnar's 1973 proof, cited in
Section~\ref{sec:clock}, that the problem is inherent to synchronizing
two unrelated clocks.

\paperref{Wikipedia, ``Flip-flop (electronics)''}{https://en.wikipedia.org/wiki/Flip-flop_(electronics)}
\end{researchnote}

With the machine's fastest memory understood bit by bit, the next
chapter returns to the programmer's view: Chapter~3 combines scalars
into \emph{structures}, and shows how the machine's demand for
aligned access quietly inserts invisible padding bytes between our
fields.
